% Generated by tools/edn_latex.py from reports/edn.md.
% Do not edit: change reports/edn.md and run `make edn`.
\ednentry{
  id = {EDN-07},
  title = {Raw data hosting: import from Zenodo, never push to our own remote},
  note = {\textbf{Amended by EDN-25} (2026-09-26, PR \#27): the raw file is tracked with \texttt{dvc add} and pushed to our DagsHub remote instead.
This entry is kept as the record of what was decided on 2026-09-22 and of the reasoning EDN-25 revised.},
  date = {2026-09-22},
  milestone = {M2: Reproducibility},
  topic = {Data Versioning, Privacy},
  participants = {@lukas2510},
  decision = {The raw AutoScout24 CSV is pulled with \texttt{dvc import-\allowbreak{}url} from its Zenodo DOI, pinned to version 1.0.0, and is never \texttt{dvc push}ed to our DagsHub remote. Only the PII-stripped \texttt{data/\allowbreak{}interim} output onward is tracked and pushed like a normal pipeline artefact. Enforced by \texttt{push: false} on the file's output in \texttt{data/\allowbreak{}raw/\allowbreak{}cars.\allowbreak{}csv.\allowbreak{}dvc}, which makes \texttt{dvc push} skip it, and a test that the field stays there.},
  rationale = {\emph{Alternatives considered:}
\begin{itemize}[leftmargin=*, topsep=2pt]
\item \textbf{Option A (chosen): \texttt{dvc import-\allowbreak{}url} from Zenodo, raw file never pushed.}
\newline Pros: the raw file carries PII (\texttt{street}, \texttt{zip}, exact coordinates, \texttt{seller\_\allowbreak{}company\_\allowbreak{}name} for private sellers); the Zenodo DOI is already immutable and versioned, so re-fetching it with \texttt{dvc update} gives the same reproducibility as a cached copy without a second, PII-bearing copy on our own remote (likely public); matches the data-minimization principle for sensitive raw data.
\newline Cons: fetching the raw data needs live access to Zenodo; reproducibility depends on Zenodo keeping this exact file available; \texttt{dvc push} uploads imported files like any other output, so the ``never pushed'' half needs \texttt{push: false} on the output, not just the command choice, to actually hold; because the file is never on the remote, a fresh clone gets it with \texttt{dvc update} instead of \texttt{dvc pull}, and \texttt{dvc repro} fails until it has.
\item \textbf{Option B: \texttt{dvc add} the raw file like any other pipeline input and push it to DagsHub.}
\newline Pros: consistent with how every other artefact in the project is tracked; does not depend on Zenodo staying reachable.
\newline Cons: re-hosts a PII-bearing file on our own remote, which is likely public; the second copy adds no reproducibility over the immutable Zenodo DOI, only exposure.
\end{itemize}
\par
\emph{Why:} Option A was already the direction noted (but left \texttt{[open]}) in the project brief's data-issues section. It was surfaced as an undecided requirement-vs-brief contradiction during the review of PR \#19 (``requirements.md'' stated the raw data is never re-hosted as a settled \texttt{NFR-\allowbreak{}08}, while the brief still called the same question open, and the existing \texttt{data-\allowbreak{}versioning.\allowbreak{}md} example showed the opposite pattern, \texttt{dvc add data/\allowbreak{}raw/\allowbreak{}cars.\allowbreak{}csv}). Data minimization for a file containing indirect identifiers of private sellers outweighs the convenience of a uniform tracking pattern; the Zenodo DOI's own immutability and versioning already give the reproducibility that pushing a copy would otherwise provide.},
  ai = {Information seeking, Alternative assessment, Recommendation},
  response = {Accepted},
  assessment = {While reviewing PR \#19 for SOTA fit and course requirements, AI noticed that \texttt{NFR-\allowbreak{}08}'s ``never re-hosted'' wording silently settled the brief's open \texttt{dvc import-\allowbreak{}url} question without going through the EDN process, and that the project's own \texttt{data-\allowbreak{}versioning.\allowbreak{}md} example contradicted it. Asked for a recommendation, AI laid out options A and B with pros and cons, recommended A on privacy and data-minimization grounds, and flagged that \texttt{dvc import-\allowbreak{}url} alone does not prevent a re-push and needs a CI check to be enforceable. Lukas reviewed the reasoning and confirmed option A. On 2026-09-23, asked whether NFR-08 was overkill, AI proposed \texttt{push: false} on the output instead of that CI check against the DagsHub remote, and verified it in a scratch repository with DVC 3.67.1: without the field \texttt{dvc push} uploads the imported file, with it the remote stays empty, also after a \texttt{dvc update} from a changed source. The same test showed that in a fresh clone \texttt{dvc repro} does not re-fetch an imported file and a plain \texttt{dvc pull} exits with an error, while \texttt{dvc update} does fetch it, which corrected the instructions in \texttt{data-\allowbreak{}versioning.\allowbreak{}md}. Lukas accepted the change of mechanism; the decision itself stayed the same.},
  aievidence = {Claude Code session on 2026-09-22 while reviewing PR \#19: AI flagged the NFR-08/brief contradiction, explained it in detail on request, then gave a recommendation and cited data-minimization and Zenodo DOI immutability as the SOTA rationale for option A; Lukas replied ``yes please do so''. Claude Code session on 2026-09-23: Lukas asked (translated from German) ``NFR-08 privacy, does that still need adjusting? how would you adjust it so it is no longer excessive'', AI proposed \texttt{push: false} with the scratch-repository evidence; Lukas replied ``yes please do that''.},
  evidence = {\href{https://github.com/mlops-2627q1-mds-upc/MLOPS_Recommenditos/blob/main/docs/docs/requirements.md}{requirements} NFR-08; \href{https://github.com/mlops-2627q1-mds-upc/MLOPS_Recommenditos/blob/main/docs/docs/specification.md}{specification} NFR-08 and ``Decisions pending confirmation''; \href{https://github.com/mlops-2627q1-mds-upc/MLOPS_Recommenditos/blob/main/docs/docs/data-versioning.md}{Data versioning}; \href{https://github.com/mlops-2627q1-mds-upc/MLOPS_Recommenditos/blob/main/docs/docs/project-brief.md}{project brief} section 3.3; \href{https://github.com/mlops-2627q1-mds-upc/MLOPS_Recommenditos/pull/19}{PR \#19}.},
}
